\chapter{格点规范理论的基础对象}

\section{链接、平行移动与 Wilson 线}

格点规范场的自由度位于有向链接上。采用本节统一的显式耦合约定
\begin{equation}
  U_\mu(x)=\mathcal P\exp\!\left[i g_0\int_x^{x+a\hat\mu}\!\dd y^\nu A_\nu(y)\right],
  \qquad
  U_{-\mu}(x)=U_\mu(x-a\hat\mu)^\dagger ,
  \label{eq:link-definition}
\end{equation}
其中 \(A_\mu=A_\mu^aT^a\)，生成元满足 \(\Tr(T^aT^b)=\tfrac12\delta^{ab}\)。链接不是普通相位，而是把颜色参考系从 \(x\) 平行移动到 \(x+a\hat\mu\) 的群元素。若把 \(g_0A_\mu\) 合并记为 \(A_\mu\)，后续所有场强与 Clover 公式都应删去显式 \(g_0\)；两种约定不能混用。

在规范变换 \(\psi(x)\to\Omega(x)\psi(x)\) 下，
\begin{equation}
  U_\mu(x)\longrightarrow \Omega(x)U_\mu(x)\Omega^\dagger(x+a\hat\mu).
  \label{eq:link-gauge}
\end{equation}
沿路径的链接乘积因此只在两个端点变换。开 Wilson 线必须与端点场或场强收缩，闭合 Wilson 圈则取颜色迹，二者才构成规范不变量。这个简单的端点规则决定了准 PDF、准 TMD 和胶子非局域算符中路径方向、表示与指标顺序的全部物理含义。

连续 Wilson 线是路径编序指数
\begin{equation}
  W_\Gamma(y,x)=\mathcal P\exp\!\left[i g_0\int_\Gamma \dd z^\mu A_\mu(z)\right].
  \label{eq:wilson-line}
\end{equation}
格点上它由沿途链接的有序乘积代替。路径不是数值实现细节：改变两条长边的方向会改变软胶子重求和，改变颜色表示会改变 Casimir 与 cusp 反常维，交换端点和负方向则改变厄米共轭关系。

\section{Wilson 圈、静态势与有限温度}

设 \(W_{rT}\) 是空间距离为 \(r\)、时间跨度为 \(T\) 的矩形 Wilson 圈，并按颜色维数归一化。取其大 \(T\) 极限，
\begin{equation}
  \langle W_{rT}\rangle
  \sim \exp[-T V(r)]\left[1+\Ord(e^{-T\Delta E})\right],
  \label{eq:wilson-loop-potential}
\end{equation}
则 \(V(r)\) 是静态夸克--反夸克势。红外区域常用
\begin{equation}
  V(r)=A+\frac{B}{r}+\sigma r
  \label{eq:static-potential}
\end{equation}
描述自能常数、Coulomb 项和线性禁闭项。面积律 \(\langle W_C\rangle\propto e^{-\sigma\mathcal A_C}\) 是禁闭的格点判据之一；强耦合展开可给出面积律的首项，但不能代替连续极限下的非微扰测量。振幅、自能和常数项依赖归一化，只有斜率与差值的连续极限才具有直接的物理意义。

静态势也连接标度设定。Sommer 参数定义为
\begin{equation}
  r_0^2\left.\frac{\dd V}{\dd r}\right|_{r=r_0}=1.65.
  \label{eq:sommer}
\end{equation}
格点先计算无量纲 \(r_0/a\)；只有选定 \(r_0\) 的物理约定后，才能赋值 \(a\)。因此“\(r_0\simeq0.5~{\rm fm}\)”是外部标度输入，不是格点理论自身的预言。

在有限温度 \(T=1/(aN_T)\) 下，沿周期时间方向的闭合 Wilson 线给出 Polyakov 圈
\begin{equation}
  P(\mathbf x)=\Tr\prod_{j=0}^{N_T-1}U_4(\mathbf x,j).
  \label{eq:polyakov}
\end{equation}
纯规范理论具有中心对称性 \(P\to zP\)，\(z\in\mathbb Z_3\)；\(\langle P\rangle=0\) 区分禁闭相，\(\langle P\rangle\ne0\) 对应中心对称破缺。动力学费米子显式破缺中心对称性，因此在全 QCD 中 Polyakov 圈不是严格有序参量，禁闭--退禁闭应作为 crossover 处理。

\section{从 plaquette 到场强张量}

基本 plaquette 定义为
\begin{equation}
  P_{\mu\nu}(x)=U_\mu(x)U_\nu(x+\hat\mu)
  U_\mu^\dagger(x+\hat\nu)U_\nu^\dagger(x).
  \label{eq:b01-plaquette}
\end{equation}
对小 \(a\) 依次使用 Baker--Campbell--Hausdorff 公式和 \(A_\nu(x+a\hat\mu)=A_\nu(x)+a\partial_\mu A_\nu(x)+\cdots\)，得到
\begin{equation}
  P_{\mu\nu}(x)
  =\exp\!\left[i a^2g_0F_{\mu\nu}(x)+\Ord(a^3)\right],
  \qquad
  F_{\mu\nu}=\partial_\mu A_\nu-\partial_\nu A_\mu
  +ig_0[A_\mu,A_\nu].
  \label{eq:plaquette-field-strength}
\end{equation}
若在此推导中把 \(g_0\) 吸收进 \(A_\mu\)，指数中应写成 \(i a^2F_{\mu\nu}\)，场强对易项也相应地不含显式 \(g_0\)。

为减小单个 plaquette 的统计噪声并提高超立方对称性，取四个定向共面 plaquette 的对称和 \(Q_{\mu\nu}\)。所有定向都应使其领头项为 \(\mathbb 1+i a^2g_0F_{\mu\nu}\)。例如在常用定向约定下可写
\begin{equation}
  Q_{\mu\nu}=P_{\mu\nu}+P_{\nu,-\mu}
  +P_{-\mu,-\nu}+P_{-\nu,\mu}.
  \label{eq:clover-sum}
\end{equation}
于是可直接定义反厄米的格点曲率
\begin{equation}
  F^{c}_{\mu\nu}(x)
  =-\frac{i}{8g_0a^2}
  \left[Q_{\mu\nu}(x)-Q_{\mu\nu}^\dagger(x)\right].
  \label{eq:clover-field}
\end{equation}
在基本表示中还应去掉颜色迹：
\begin{equation}
  F_{\mu\nu}\longmapsto
  F_{\mu\nu}-\frac1{N_c}\Tr(F_{\mu\nu})\,\mathbb 1 .
  \label{eq:traceless-field}
\end{equation}
这一步应发生在场强与 Wilson 线收缩之前。式~\eqref{eq:clover-field} 的 \(g_0\) 与 \(a^2\) 因子是和式~\eqref{eq:link-definition} 绑定的一整套约定；若改用 \(U=e^{iaA}\) 或不同 \(Q\) 归一化，必须重新推导而不是只替换一个符号。

格点场强具有 \(\Ord(a^2)\) 的离散化误差。固定物理尺度并减小 \(a\) 时，它趋于连续场强；靠近截断时，单 plaquette 的高维算符、有限体积镜像和与平滑链接不一致的端点都会进入系统误差。对胶子算符而言，场强、Wilson 线、迹投影和 Fourier 约定属于同一个算符定义，不能分别“调参”。

\chapter{离散化误差、改进与标度设定}

\section{Symanzik 有效展开}

格点关联函数可在长距尺度上用连续 QCD 加局域高维算符展开：
\begin{equation}
  S_{\rm eff}=S_{\rm QCD}
  +a\int\!\dd^4x\,\mathcal L_1(x)
  +a^2\int\!\dd^4x\,\mathcal L_2(x)+\cdots .
  \label{eq:symanzik}
\end{equation}
\(\mathcal L_n\) 由连续场、连续导数和格点对称性允许的算符组成，量纲为 \(4+n\)。该展开说明离散误差可按 \(a\) 幂次系统分类；它不意味着在有限 \(a\) 上删除了 UV 自由度，也不意味着所有高阶系数都能由树级计算得到。

标准 Wilson 规范作用量从 \(\Ord(a^2)\) 偏差开始。树级 Symanzik 改进在 plaquette 之外加入 \(1\times2\) 矩形圈。在独立平面、同一方向归一化的常用约定中，树级系数为
\begin{equation}
  c_0=\frac53,\qquad c_1=-\frac1{12},
  \qquad c_0+8c_1=1.
  \label{eq:symanzik-coefficients}
\end{equation}
最后一个等式是检验系数与求和约定是否自洽的最简单条件；改变矩形圈的计数方式会改变组合因子，不能只引用 \(-1/12\) 而省略求和定义。

Wilson 费米子的主导误差为 \(\Ord(a)\)。Sheikholeslami--Wohlert 项
\begin{equation}
  S_f=S_W+a^5c_{\rm SW}\sum_x
  \bar\psi(x)\frac{i}{4}\sigma_{\mu\nu}F_{\mu\nu}(x)\psi(x)
  \label{eq:b01-clover-action}
\end{equation}
在适当调谐 \(c_{\rm SW}\) 后消除物理矩阵元中的 \(\Ord(a)\) 项。树级 \(c_{\rm SW}=1\)，单圈形式为 \(1+\Ord(g_0^2)\)，非微扰值可由 Schrödinger 泛函或等价的手征 Ward 条件确定。“使用 Clover”本身不是精度保证：残余 \(\Ord(a\alpha_s)\)、重夸克 \(\Ord(a^2m_Q^2)\)、算符混合和不完整的手征外推仍须逐项估计。

\section{Tadpole 改进与作用量一致性}

格点微扰论中的蝌蚪图使链接的真空期望值明显小于单位矩阵。一个常用非微扰量是
\begin{equation}
  u_0=\left\langle\frac1{N_c}\Re\Tr P_{\mu\nu}\right\rangle^{1/4},
  \qquad U_\mu\longmapsto U_\mu/u_0 .
  \label{eq:tadpole}
\end{equation}
重标链接后，作用量、Clover 系数、临界质量与算符归一化必须同步改变。例如在本节 Clover 归一化下，未重标链接的树级值为 \(c_{\rm SW}=1\)；显式采用 \(U/u_0\) 后，常用树级值为 \(u_0^{-3}\)，临界 hopping 参数也要按相应约定换算。把两种表述混在一起会制造伪装的 \(\Ord(a)\) 误差。

规范作用和费米子作用若都采用 tadpole 改进，应使用同一套自洽定义。不同 \(u_0\) 或只在费米子部分重标，会使规范场和费米子场在不同的有效归一化上“运作”。树级改进之后仍可能存在 \(\Ord(a\alpha_s)\) 交叉项；多格距拟合若只保留 \(a^2\)，该误差可能被错误吸收到物理常数中。

\section{标度设定不是物理极限}

格点计算首先给出无量纲比值。固定一个物理参考长度 \(\ell_{\rm ref}\) 后，
\begin{equation}
  a=\frac{\ell_{\rm ref}}{\ell_{\rm ref}/a}.
  \label{eq:b01-scale-setting}
\end{equation}
\(r_0\)、梯度流尺度 \(w_0\) 或强子质量都可充当 \(\ell_{\rm ref}\)，但它们的统计噪声、夸克质量依赖和外部输入不同。梯度流的常用定义是
\begin{equation}
  \frac{\dd B_\mu(x,t)}{\dd t}
  =-g_0^2\frac{\delta S_W[B]}{\delta B_\mu(x,t)},
  \qquad
  t^2\langle E(t)\rangle\big|_{t=t_0}=0.3,
  \label{eq:flow-scale}
\end{equation}
其中 \(w_0^2=t_0\) 是由同一条件定义的两种记法。流半径 \(r_f=\sqrt{8t}\) 控制平滑尺度；它不是一个可以任意选择后忽略的“数值参数”。

在固定格距族上，改进后的物理量常写成
\begin{equation}
  X(a,m_\pi)=X_0+d_a a^2+d_4a^4
  +d_\pi(m_\pi^2-m_{\pi,\rm phys}^2)+\cdots .
  \label{eq:continuum-fit}
\end{equation}
是否加入 \(a^4\)、交叉项 \(a^2P_z^2\) 或 \(a^2m_Q^2\) 取决于目标和格距覆盖，而不能由数据的 \(\chi^2\) 单独决定。至少两个格距只能检查一个格距趋势；要对多项同时拟合，需要足够多的独立格距和可控的相关误差。

\chapter{费米子谱、色散关系与组态采样}

\section{自由格点色散与加倍子}

最近邻自由标量场的传播子在 Euclidean 动量空间为
\begin{equation}
  G^{-1}(p)=\frac4{a^2}\sum_{\mu=1}^4\sin^2\!\left(\frac{ap_\mu}{2}\right)+m^2.
  \label{eq:scalar-propagator}
\end{equation}
把 \(p_4=iE(\mathbf p)\) 代入 \(G^{-1}=0\)，得到
\begin{equation}
  \cosh[aE(\mathbf p)]
  =\cosh(am)+\sum_{i=1}^3\left[1-\cos(ap_i)\right].
  \label{eq:lattice-dispersion}
\end{equation}
在固定物理动量并令 \(a\to0\) 时，
\begin{equation}
  E^2(\mathbf p)=m^2+\mathbf p^2+\Ord(a^2).
  \label{eq:continuum-dispersion}
\end{equation}
四维超立方对称性通常排除一般 \(\Ord(a)\) 项；若写 \(\Ord(ap)\)，应把它理解为对可用动量范围过宽的保守估计，或明确成 \(\Ord(a^2E^2)\)。在有限格距上，式~\eqref{eq:lattice-dispersion} 可用于拟合高 boost 强子；若数据偏离连续色散，必须先区分激发态、动量涂抹和格点色散，而不能直接归因于新物理。

naive Dirac 算符在动量空间含有 \(\sin(ap_\mu)\)。因此每个动量分量在 \(0\) 与 \(\pi/a\) 附近都可能贡献低能极点，四维共出现 \(2^4=16\) 个候选模。Wilson 项
\begin{equation}
  D_W(p)=m+\frac{i}{a}\sum_\mu\gamma_\mu\sin(ap_\mu)
  +\frac{r}{a}\sum_\mu\left[1-\cos(ap_\mu)\right]
  \label{eq:wilson-dirac}
\end{equation}
给布里渊区边界附近的模增加正比于 \(1/a\) 的有效质量，同时显式破坏标准手征对称性。对空间零动量、正能分支，自由场关系常可化为 \(e^{aE}=1+am\) 的特定形式；它只适用于该分支，不能当作任意动量的精确色散。

Nielsen--Ninomiya 约束表明：平移不变、局域、厄米并保持标准手征反对易关系的格点费米子不能同时避免加倍子。实际方案必须改变其中一个条件。Wilson/Clover 显式破缺手征对称性；Staggered/HISQ 保留 \(U(1)\) 子群但引入 taste；Domain-Wall/Overlap 利用第五维或 Ginsparg--Wilson 非局域性；Twisted Mass 通过最大扭曲实现自动 \(\Ord(a)\) 改进。表~\ref{tab:fermions} 给出选择时应同时检查的性质。

{\small
\begin{longtable}{L{2.4cm}L{1.8cm}L{2.1cm}L{2.4cm}L{2.8cm}}
\caption{主要格点费米子方案的结构性比较。表中“成本”是相对方法级判断，不是具体机器性能测量。}\label{tab:fermions}\\
\toprule
方案 & 手征性 & 主导离散误差 & 主要代价/风险 & 典型用途\\
\midrule
\endfirsthead
\toprule
方案 & 手征性 & 主导离散误差 & 主要代价/风险 & 典型用途\\
\midrule
\endhead
Wilson & 显式破缺 & \(\Ord(a)\) & 加法质量重整化、例外组态 & 基础谱学与大质量夸克\\
Clover & 显式破缺 & \(\Ord(a^2)\)（调谐后） & \(c_{\rm SW}\)、\(a\alpha_s\) 与算符混合 & 强子结构与部分子算符\\
Staggered\slash HISQ & \(U(1)\) 子群 & \(\Ord(a^2)\)，HISQ 进一步压低 & taste 分裂、根技巧与粗格距限制 & 大体积动力学系综\\
Domain-Wall & 有限 \(L_s\) 近似 & \(m_{\rm res}\) 控制后 \(\Ord(a^2)\) & 第五维成本、残余质量 & 手征极限与交叉检验\\
Twisted Mass & 显式破缺 & 最大扭曲下 \(\Ord(a^2)\) & 同位旋破缺、临界质量调谐 & 轻海夸克与谱学\\
Overlap & 精确 GW & \(\Ord(a^2)\) & 矩阵符号函数的成本、拓扑冻结 & 手征对称性基准与验证\\
\bottomrule
\end{longtable}
}

\section{Markov 链与 Metropolis 基础}

格点路径积分的采样目标可写成
\begin{equation}
  P(U)=\frac{e^{-S_g[U]}\det D[U]}{\int\mathcal DU'\,e^{-S_g[U']}\det D[U']}.
  \label{eq:target-weight}
\end{equation}
对 Markov 转移 \(T(U'|U)\)，细致平衡条件
\begin{equation}
  P(U)T(U'|U)=P(U')T(U|U')
  \label{eq:detailed-balance}
\end{equation}
保证 \(P\) 是平稳分布。对称提议的 Metropolis 接受概率为
\begin{equation}
  T_{\rm acc}(U'|U)=\min\!\left[1,e^{-(S[U']-S[U])}\right].
  \label{eq:metropolis}
\end{equation}
单链接扫描只改变与链接相邻的 plaquette；热浴直接从局部条件分布抽样，过松弛则用保作用量的确定性移动加速遍历。实际误差不能只写 \(\sigma/\sqrt N\)：链样本存在自相关，应报告积分自相关时间和有效样本数。

\section{赝费米子与 HMC}

对两味简并费米子，使用正定矩阵 \(D^\dagger D\)：
\begin{equation}
  \int\mathcal D\phi^\dagger\mathcal D\phi\,
  e^{-\phi^\dagger(D^\dagger D)^{-1}\phi}
  =\det(D^\dagger D).
  \label{eq:b01-pseudofermion}
\end{equation}
抽取 \(\chi\sim\exp(-\chi^\dagger\chi)\) 并令 \(\phi=D\chi\)，即可得到协方差为 \(D^\dagger D\) 的高斯场。单味或非偶味情形需要用有理/多项式近似（RHMC/PHMC），并分别控制近似阶数与力误差。

HMC 扩展相空间为 \((U,\Pi)\)。取
\begin{equation}
  H[U,\Pi]=\frac12\Tr\Pi^2+S_g[U]
  +\phi^\dagger(D^\dagger D)^{-1}\phi,
  \label{eq:hmc-hamiltonian}
\end{equation}
通过蛙跳积分得到候选终态。定义 \(F=-\delta S/\delta\omega\) 后，可写为
\begin{equation}
  \Pi_{1/2}=\Pi_0-\frac{\epsilon}{2}F[U_0],\qquad
  U_{k+1}=e^{i\epsilon\Pi_{k+1/2}}U_k,
  \qquad
  \Pi_{k+3/2}=\Pi_{k+1/2}-\epsilon F[U_{k+1}].
  \label{eq:leapfrog}
\end{equation}
末端以 \(\min(1,e^{-\Delta H})\) 接受。精确算术下蛙跳可逆并保持相空间体积，有限步长的能量偏差为 \(\Ord(\epsilon^2)\)，由全局接受步骤修正。

HMC 的生产质量至少需要报告：接受率、轨迹长度、步长、力、Dirac 真残差、拓扑荷历史、不同拓扑扇区的采样和自相关时间。拓扑冻结、近零模和力发散不能由“接受率看起来正常”排除。HMC 的产物是组态集合，不是物理矩阵元；后续两点/三点函数、谱学拟合和重正化仍各有独立误差。

\chapter{非定域算符与重正化}

\section{三类 UV 结构不能混为一谈}

含 Wilson 线的裸矩阵元通常表现为
\begin{equation}
  h_{\rm bare}(z,b_T,L;a)
  \sim Z_{\rm endpoint}(a,\mu)\,
  e^{-\delta m L/a}\,
  e^{-V(b_T)L}\,
  h_{\rm phys}(z,b_T;\mu),
  \label{eq:nonlocal-structure}
\end{equation}
其中第一项是端点对数与算符反常维，第二项是 Wilson 线自能的线性项，第三项是两条近平行长 Wilson 线之间交换胶子的 pinch-pole 结构。若 Wilson 线沿光锥方向延伸到无穷，还会出现 rapidity 发散；rapidity 不是普通 UV 发散，不能用 \(\delta m\) 或局部反常量纲替代。

\begin{longtable}{L{2.3cm}L{3.4cm}L{2.8cm}L{4.0cm}}
\caption{非定域算符中常见发散及所需控制方式。}\label{tab:b01-divergences}\\
\toprule
结构 & 物理来源 & 典型尺度 & 控制方式\\
\midrule
\endfirsthead
\toprule
结构 & 物理来源 & 典型尺度 & 控制方式\\
\midrule
\endhead
线性自能 & 开 Wilson 线的静态色源自能 & \(L/a\) 或 \(|z|/a\) & 质量抵消、同长度 Wilson 圈平方根、多格距自重整\\
端点对数 & 场强/夸克顶点与 Wilson 线连接 & \(\ln(1/a)\) & 反常量纲、RI/MOM 端点因子、短距匹配\\
pinch pole & 两条长 Wilson 线之间的有效势 & \(V(b_T)L\) & 同几何 Wilson 圈、\(L\to\infty\) 平台\\
cusp & 路径转角与有限夹角 & 转角反常维 & 路径匹配的 Wilson 圈减除\\
rapidity & 光锥方向软/共线积分 & \(L_\parallel\) 或 \(\zeta\) & soft function、rapidity regulator、CS 演化\\
\bottomrule
\end{longtable}

\section{局部 RI/MOM 与它的边界}

对局域夸克双线性算符 \(O_\Gamma=\bar u\Gamma d\)，RI/MOM 条件可概括为
\begin{equation}
  Z_\Gamma\frac1{12}\Tr\!\left[
  \langle p|O_\Gamma|p\rangle
  \langle p|O_\Gamma|p\rangle_0^{-1}
  \right]_{p^2=\mu^2}=1,
  \label{eq:ri-mom}
  \end{equation}
其中下标 \(0\) 表示树级或自由场归一化，\(\mu\) 必须落在微扰窗口
\begin{equation}
  \Lambda_{\rm QCD}^2\ll\mu^2\ll a^{-2}.
  \label{eq:perturbative-window}
\end{equation}
RI/MOM 可以确定局部反常量纲和有限混合矩阵，并通过微扰转换到 \(\MSbar\)。但离壳顶点条件不知道开 Wilson 线的总长度与同几何自能；把普通 RI/MOM 直接用于长 Wilson 线或 staple，通常会留下线性项。因此，RI/MOM 失效不表示所有动量减除都无用，而是说明必须补上非定域几何对应的减除。

\section{质量抵消、Wilson 圈与自重整}

辅助场方法把开 Wilson 线看成无动能静态色源的两点函数。若辅助场 \(Z\) 具有运动方程
\begin{equation}
  [\partial_z-ig_0A_z(z)]Z(z)=0,
  \qquad L(z,0)=Z(z)Z^\dagger(0),
  \label{eq:auxiliary-z}
\end{equation}
其线性自能可由质量抵消项吸收。对固定几何和表示，可把重整化结构写成
\begin{equation}
  L_R(z,0)=Z_Z^{-1}e^{-\delta m|z|}L(z,0),
  \label{eq:wilson-renorm}
\end{equation}
其中 \(Z_Z\) 仍包含端点对数。

对 staple 型准 TMD，同几何 Wilson 圈减除写成
\begin{equation}
  h_R(z,b_T,L;a)
  =\lim_{L\to\infty}
  \frac{h_{\rm bare}(z,b_T,L;a)}
  {\sqrt{Z_E(2L+z,b_T;a)}} ,
  \label{eq:wilson-loop-subtraction}
\end{equation}
其中 \(Z_E\) 的矩形圈必须与开 Wilson 线具有匹配的总长度、转角、横向分离和颜色表示。若线性项与 pinch 项都按长度成比例，开方的圈可抵消它们；残余端点对数再通过短距比率（SDR）或混合方案匹配到 \(\MSbar\)。

自重整化利用同一矩阵元的比值
\begin{equation}
  h_R(z,P_z)
  =\frac{h_{\rm bare}(z,P_z)}
  {\left[h_{\rm bare}(z_0,P_z=0)\right]^{z/z_0}},
  \label{eq:self-renorm}
\end{equation}
以短距点 \(z_0\) 标定共同指数因子。当 $z/z_0$ 为整数或通过内插定义时，它可消除与 \(z\) 成正比的共同线性项；但它是否同时消除端点对数、离散项和 rapidity 结构，必须用多格距稳定性与方案匹配证明。单格距下 \(z_0\) 与参数化模型往往不可辨识。

混合方案把短距与长距分开：
\begin{equation}
  h_R(z,P_z)=
  \begin{cases}
    h_{\rm bare}(z,P_z)/h_{\rm bare}(z,0), & z\le z_s,\\[2mm]
    h_{\rm bare}(z,P_z)/
    \left[Z_R(z,a,\mu)h_{\MSbar}(z_s,0,\mu)\right], & z>z_s,
  \end{cases}
  \label{eq:hybrid}
\end{equation}
切换点必须满足 \(a\ll z_s\ll\Lambda_{\rm QCD}^{-1}\)，并检查两侧在 \(z_s\) 处是否连续、最终结果是否依赖 \(z_s\)。若切换点变化超过目标误差，幂修正或微扰匹配尚未受控。

\section{梯度流与重整子}

梯度流将规范场向平滑配置演化。在固定流时下，规范不变的流场关联函数可以具有受控连续极限；其常用尺度窗口为
\begin{equation}
  a\ll r_f=\sqrt{8t}\ll
  \min\left(z,b_T,\Lambda_{\rm QCD}^{-1}\right).
  \label{eq:flow-window}
\end{equation}
左端保证 UV 平滑尺度大于格距，右端保证流没有抹平待测长距结构。流化算符通过短流时展开与普通重整化算符联系：
\begin{equation}
  O(t)=c_O(t,\mu)O_R(\mu)
  +\sum_i c_i(t,\mu)O_{i,R}(\mu)+\cdots ,
  \label{eq:flow-matching}
\end{equation}
因此流时不是普通 UV 标度，也不能替代 rapidity 重正化或算符混合矩阵。

在微扰展开中，Wilson 系数往往是渐近级数。红外重整子使大阶系数按
\begin{equation}
  c_n\sim n!\left(\frac{\beta_0}{2\pi}\right)^n
  \label{eq:renormalon-growth}
\end{equation}
增长，最小项误差对应更高 twist 的幂次模糊。若用主值（PV）或其他 Borel 重求和定义领头重整子，则必须把相应的 \(m_0z\) 线性项、短距 Wilson 系数和匹配阶数放在同一方案中。只对一个局部分量做重求和、却保持其他分量在另一方案，会人为制造方案差。

\chapter{有效场论连接：Symanzik、HQET 与 LaMET}

三种展开的共同骨架是“大标度 + 微扰匹配 + 幂次修正”，但大标度不同：
\begin{center}
\begin{tabularx}{\textwidth}{L{2.4cm}L{3.2cm}L{3.2cm}Y}
\toprule
框架 & 大标度 & 小参数 & 主要失效边界\\
\midrule
Symanzik & \(a^{-1}\) & \(a\Lambda_{\rm QCD}\) & 靠近截断、高阶改进系数和非局域长程算符\\
HQET & \(m_Q\) & \(\Lambda_{\rm QCD}/m_Q\) & 粲夸克、\(am_Q\) 非小、强末态与有限质量修正\\
LaMET & \(P_z\) & \(\Lambda_{\rm QCD}/P_z\) & 端点、大 \(z\)、格点色散和 \(aP_z\) 非小\\
\bottomrule
\end{tabularx}
\end{center}

HQET 从 QCD 场的速度分解开始：
\begin{equation}
  Q(x)=e^{-im_Qv\cdot x}\left[h_v(x)+H_v(x)\right],
  \qquad v^2=1,
  \label{eq:hqet-decomposition}
\end{equation}
其中 \(v\!\!/h_v=h_v\)，\(v\!\!/H_v=-H_v\)。积掉小分量后，
\begin{equation}
  \mathcal L_{\rm HQET}
  =\bar h_viv\cdot D h_v
  +\frac{1}{2m_Q}\bar h_v(iD_\perp)^2h_v
  +\frac{g_0}{4m_Q}\bar h_v\sigma_{\mu\nu}G^{\mu\nu}h_v
  +\Ord(m_Q^{-2}).
  \label{eq:hqet-lagrangian}
\end{equation}
领头项显示重夸克自旋--味对称性；动能和色磁项分别描述残余运动与超精细分裂。HQET 的匹配因子吸收高标度 UV 物理，非微扰矩阵元包含低能强子结构。对格点重味，必须同时报告 \(m_Qa\)、\(\Lambda_{\rm QCD}a\) 和重整化标度；只做形式上的 \(1/m_Q\) 展开而不控制 \(am_Q\)，不能得到可控连续极限。

LaMET 可以写成类似的有效展开：
\begin{equation}
  F(P_z/\Lambda)
  =Z\!\left(\frac{P_z}{\Lambda},\frac{\Lambda}{\mu}\right)f(\mu)
  +\Ord\!\left(\frac{\Lambda^2}{P_z^2},\frac{\Lambda}{P_z}\right).
  \label{eq:lamet-expansion}
\end{equation}
最低阶的类比很有用：\(P_z\leftrightarrow m_Q\)，光锥分布 \(f\leftrightarrow\) HQET 低能矩阵元。但 LaMET 的算符还带 Wilson 线几何、rapidity 和软函数，不能仅凭 HQET 类比推出胶子 TMD 的因子化。对重介子 LCDA，若
\begin{equation}
  P_z\gg m_H\gg\Lambda_{\rm QCD},
  \label{eq:bhqet-hierarchy}
\end{equation}
则可先匹配准分布到 QCD LCDA，再匹配到 boosted HQET LCDA。两个匹配核和两个幂修正参数必须独立记录；把 \(m_H/P_z\) 与 \(\Lambda_{\rm QCD}/m_H\) 合并成一个“小参数”会失去可检验性。

\chapter{综合验收流程}

从格点组态到物理结论，建议按下列顺序设置门控：
\begin{enumerate}
  \item \textbf{定义门：}固定链接方向、生成元、场强符号、Wilson 线表示、Fourier 约定、态归一化和算符端点。
  \item \textbf{代数门：}在恒等链接、纯规范变换和小矩阵自由场中检查协变性、厄米性、Clover 归一化和色迹投影。
  \item \textbf{采样门：}报告 HMC/RHMC 参数、接受率、力、真残差、拓扑荷与自相关；没有有效样本数就不报告统计误差。
  \item \textbf{谱学门：}用至少两个源/汇结构和多个动量检查基态平台、色散关系与激发态污染。
  \item \textbf{重正化门：}区分局部反常量纲、线性 Wilson 线自能、pinch/cusp、rapidity；验证同几何因子、短距窗口和多格距稳定性。
  \item \textbf{连续极限门：}分别扫描 \(a\)、\(P_z\)、\(L\)、流时和匹配标度；任何极限顺序变化都必须成为系统误差项。
  \item \textbf{物理论述门：}只有所有门均通过，才把矩阵元称为目标 PDF、TMD 或重味物理量；否则只能报告诊断量或中间矩阵元。
\end{enumerate}
